% ======================================================================
%  Приложения (обязательные): три чертежа -- общий вид, сборочный, подшипник.
%  Буквы приложений по ГОСТ -- кириллица: А, Б, В (без «С»).
%  Для готового файла рисунок вставляется автоматически; если файла ещё
%  нет в папке figures/ -- показывается рамка-плейсхолдер.
% ======================================================================

% --- Приложение А: общий вид ---
\newpage
\phantomsection
\addcontentsline{toc}{section}{Приложение А}
\begin{center}
{\normalsize\bfseries ПРИЛОЖЕНИЕ А}\\[2pt]
(обязательное)\\[\baselineskip]
{\normalsize\bfseries Общий вид насоса ЦНС 105-392}
\end{center}

\vspace{\baselineskip}

\begin{figure}[H]
\centering
\maybeincludegraphics{figures/fig_1_1_cns_general.png}{0.95\textwidth}{Общий вид насоса ЦНС 105-392}
\end{figure}

% --- Приложение Б: сборочный чертёж ---
\newpage
\phantomsection
\addcontentsline{toc}{section}{Приложение Б}
\begin{center}
{\normalsize\bfseries ПРИЛОЖЕНИЕ Б}\\[2pt]
(обязательное)\\[\baselineskip]
{\normalsize\bfseries Сборочный чертёж насоса ЦНС 105-392}
\end{center}

\vspace{\baselineskip}

\begin{figure}[H]
\centering
\maybeincludegraphics{figures/appendix_assembly.png}{0.95\textwidth}{Сборочный чертёж насоса ЦНС 105-392}
\end{figure}

% --- Приложение В: подшипник ---
\newpage
\phantomsection
\addcontentsline{toc}{section}{Приложение В}
\begin{center}
{\normalsize\bfseries ПРИЛОЖЕНИЕ В}\\[2pt]
(обязательное)\\[\baselineskip]
{\normalsize\bfseries Опорный подшипник скольжения}
\end{center}

\vspace{\baselineskip}

\begin{figure}[H]
\centering
\maybeincludegraphics{figures/appendix_bearing.png}{0.95\textwidth}{Опорный подшипник скольжения}
\end{figure}
